% A+ reusable source. Not yet compiled; C+ owns real compilation/rendering.
% Units use SI conventions; radians dimensionless. Requires amsmath, amssymb, amsthm.
\subsection*{Notation}
\begin{tabular}{llll}
Symbol & Meaning & Unit & Domain \\
$p$ & spiral pitch & m & p>0 \\
$b$ & radial coefficient & m & b=p/(2pi)>0 \\
$t$ & elapsed time & s & 0<=t<=T \\
$T$ & last requested time & s & T=300 \\
$v_0$ & head speed & m/s & v0=1>0 \\
$N$ & number of handles & 1 & N=224 positive integer \\
$i$ & handle index & 1 & i in\{0,...,N-1\} \\
$\theta_i(t)$ & handle angle & 1 & theta\_i>0; radians treated dimensionless \\
$\theta_{\mathrm{init}}$ & head initial angle & 1 & 32pi \\
$\mathbf X(\theta)$ & curve position & m & positive theta to R^2 \\
$\mathbf X_i(t)$ & handle position & m & X(theta\_i(t)) \\
$\boldsymbol\tau(\theta)$ & derivative of position with respect to dimensionless angle & m & theta>0 \\
$S(\theta)$ & arc primitive with S(0)=0 & m & theta>=0 \\
$u$ & dummy angle of integration & 1 & u>=0 \\
$L_i$ & straight distance for rod i & m & i=1..223; L1=2.86,others1.65 \\
$\ell$ & positive test rod distance & m & ell>0 \\
$r$ & previous-handle radius in local chord solve & m & r=b theta\_(i-1)>0 \\
$\delta$ & outward angle difference & 1 & 0<delta<pi \\
$F(\delta;r,\ell)$ & squared chord residual & m^2 & delta in[0,pi] \\
$\mathbf d_i$ & rod vector Xi-Xi-1 & m & i=1..223 \\
$\omega_i$ & angular velocity & 1/s & d theta\_i/dt \\
$\mathbf V_i$ & velocity vector & m/s & d Xi/dt \\
$U_i$ & speed magnitude & m/s & norm Vi>=0 \\
$D_i$ & velocity recurrence denominator & m^2 & d\_i dot tau(theta\_i)>0 on certified branch \\
\end{tabular}
\subsection*{Assumptions}
\paragraph{asm:rigid} Fixed straight chord lengths, ideal hinges, no hole clearance dynamics.
\paragraph{asm:branch} Follow the outward first root in (0,pi), when the sign bracket is certified.
\paragraph{asm:scope} Only handle-centre kinematics for Q1; no collisions, bench swept volume, turn-around or optimization.
\subsection*{Equations}
% Positive Archimedean spiral
\begin{equation}
b=\frac{p}{2\pi},\qquad \mathbf X(\theta)=b\theta(\cos\theta,\sin\theta)^{\mathsf T}
\label{eq:spiral}
\end{equation}
% Derivative and arc density
\begin{equation}
\boldsymbol\tau(\theta)=\mathbf X^{\prime}(\theta)=b(\cos\theta-\theta\sin\theta,\sin\theta+\theta\cos\theta)^{\mathsf T},\quad \|\boldsymbol\tau(\theta)\|=b\sqrt{1+\theta^2}
\label{eq:tangent}
\end{equation}
% Exact arc primitive, S(0)=0
\begin{equation}
S(\theta)=\int_0^\theta b\sqrt{1+u^2}\,du=\frac b2\left(\theta\sqrt{1+\theta^2}+\operatorname{arsinh}\theta\right)
\label{eq:arc}
\end{equation}
% Head arc equation and validity domain
\begin{equation}
\boxed{S(\theta_{\mathrm{init}})-S(\theta_0(t))=v_0t},\quad \theta_{\mathrm{init}}=32\pi,\quad 0\le t<S(\theta_{\mathrm{init}})/v_0
\label{eq:head}
\end{equation}
% Equivalent head ODE, sign fixed by inward motion
\begin{equation}
\boxed{\omega_0(t)=-\frac{v_0}{b\sqrt{1+\theta_0(t)^2}}}
\label{eq:head-ode}
\end{equation}
% Straight rods, not arc distances
\begin{equation}
\|\mathbf X_i(t)-\mathbf X_{i-1}(t)\|^2=L_i^2,\quad i=1,\ldots,N-1
\label{eq:rod}
\end{equation}
% Stable squared-chord residual
\begin{equation}
\boxed{F(\delta;r,\ell)=(b\delta)^2+4r(r+b\delta)\sin^2(\delta/2)-\ell^2=0}
\label{eq:stable-chord}
\end{equation}
% Strict positivity only when b,r>0 and 0<delta<pi
\begin{equation}
\frac{\partial F}{\partial\delta}=2b^2\delta+4rb\sin^2(\delta/2)+2r(r+b\delta)\sin\delta>0
\label{eq:chord-derivative}
\end{equation}
% Existence needs endpoint sign condition as well as monotonicity
\begin{equation}
F(0;r,\ell)=-\ell^2<0,\quad F(\pi;r,\ell)=(2r+b\pi)^2-\ell^2>0\quad\text{if }0<\ell<2r+b\pi
\label{eq:bracket}
\end{equation}
% Differentiate constant squared rod length; two-times factor cancels
\begin{equation}
\mathbf d_i=\mathbf X_i-\mathbf X_{i-1},\qquad \mathbf d_i^{\mathsf T}\bigl(\boldsymbol\tau(\theta_i)\omega_i-\boldsymbol\tau(\theta_{i-1})\omega_{i-1}\bigr)=0
\label{eq:velocity-constraint}
\end{equation}
% Recurrence and positive denominator on certified branch; values have different units as declared in symbol table
\begin{equation}
\boxed{\omega_i=\frac{\mathbf d_i^{\mathsf T}\boldsymbol\tau(\theta_{i-1})}{D_i}\omega_{i-1}},\quad D_i=\mathbf d_i^{\mathsf T}\boldsymbol\tau(\theta_i)=\tfrac12\frac{\partial F}{\partial\delta}>0,\quad \mathbf V_i=\boldsymbol\tau(\theta_i)\omega_i,\quad U_i=\|\mathbf V_i\|
\label{eq:velocity}
\end{equation}
% With frozen p=.55m, implies head theta>16pi until300s
\begin{equation}
S(32\pi)-S(16\pi)\ge\frac b2\left((32\pi)^2-(16\pi)^2\right)-16\pi b=105.6\pi-4.4>300\,\mathrm m
\label{eq:radius-bound}
\end{equation}
\paragraph{derive:arc}
Conditions: b>0; theta>=0.
\begin{enumerate}
\item Differentiate X componentwise using the product rule.
\item Expand squared tangent norm: cross terms cancel and cos²+sin²=1, obtaining b²(1+theta²).
\item Use d[theta sqrt(1+theta²)]/dtheta=sqrt(1+theta²)+theta²/sqrt(1+theta²) and d asinh(theta)/dtheta=1/sqrt(1+theta²).
\item The sum times b/2 is b sqrt(1+theta²); value at zero vanishes.
\end{enumerate}
\paragraph{derive:head}
Conditions: b>0; v0>0; 0<=t<S(theta\_init)/v0.
\begin{enumerate}
\item Constant speed along the curve makes traversed length v0*t.
\item Inward travel subtracts arc length from S(theta\_init).
\item Differentiate using the chain rule: -Sprime(theta0)*omega0=v0; divide by positive b sqrt(1+theta0²).
\end{enumerate}
\paragraph{derive:chord}
Conditions: r>0; b>0; ell>0.
\begin{enumerate}
\item Rotate previous handle to angle0 without changing Euclidean distance. Radii are r and r+b*delta.
\item Squared distance equals r²+(r+b*delta)²-2r(r+b*delta)cos(delta).
\item Regroup as (b*delta)²+2r(r+b*delta)(1-cos(delta)).
\item Use 1-cos(delta)=2sin²(delta/2); subtract ell². This avoids subtracting nearly equal O(r²) terms.
\end{enumerate}
\paragraph{derive:velocity}
Conditions: constant\_chords; positive\_geometry; branch\_scope; differentiable branch.
\begin{enumerate}
\item Let d\_i=X\_i-X\_i-1. Differentiate d\_i dot d\_i=L\_i²: 2d\_i dot(V\_i-V\_i-1)=0.
\item Apply chain rule V\_i=tau(theta\_i)*omega\_i; rearrange coefficients of omega\_i.
\item With previous handle fixed, derivative of squared distance with respect to delta is 2d\_i dot tau(theta\_i).
\item Use strict Fdelta>0 on the certified interval; division is justified analytically. A finite-precision threshold is still necessary for conditioning.
\end{enumerate}
\paragraph{prop:head-unique}
Conditions: b>0; v0>0; theta\_init>0; 0<=t<S(theta\_init)/v0.
\begin{enumerate}
\item S is continuous and strictly increasing because Sprime=b sqrt(1+theta²)>0.
\item Right side S(theta\_init)-v0t lies in(0,S(theta\_init)] so a unique head angle in(0,theta\_init] exists by the intermediate value theorem.
\item Differentiation gives omega0<0; t=0 returns theta\_init. The positive-radius domain ends at the origin, not beyond.
\end{enumerate}
\paragraph{prop:chord-first-root}
Conditions: positive\_geometry; constant\_chords; branch\_scope; 0<ell<2r+b*pi.
\begin{enumerate}
\item F is continuous on[0,pi]. Its lower endpoint is strictly negative and upper endpoint strictly positive under the stated chord bound.
\item On(0,pi), every term of Fdelta is nonnegative and 2b²delta is strictly positive. Thus F is strictly increasing.
\item The intermediate value theorem gives exactly one root delta* in(0,pi). No smaller positive root exists by monotonicity, so it is the first outward root.
\item No claim is made about branches outside pi or that no roots exist if the upper sign test fails.
\end{enumerate}
\paragraph{prop:frozen-interval}
Conditions: p=.55m; v0=1m/s; theta\_init=32pi; 0<=t<=300s; L\_i<=2.86m.
\begin{enumerate}
\item For u>=0, u<=sqrt(1+u²)<=u+1; integrate to bound S(32pi) from below and S(16pi) from above.
\item The resulting difference is105.6pi-4.4>300m (already follows from pi>3). Thus at300s the head theta remains>16pi and radius>4.4m; earlier radii are larger.
\item Induct along the outward chain: every previous radius is>4.4m. Hence 2r+b*pi>9.075m>2.86m, certifying each neighbour bracket throughout the continuous interval.
\item At each finite stage the implicit-function derivative Fdelta is positive, so differentiability and the velocity recurrence hold locally; no uniform floating-point condition-number bound is asserted.
\end{enumerate}
